\documentclass[11pt]{amsart}
\usepackage[utf8]{inputenc}
\usepackage[T1]{fontenc}
\usepackage{lmodern}
\usepackage{amsmath,amssymb,amsthm}
\usepackage{booktabs}
\usepackage[margin=1.1in]{geometry}
\usepackage[hidelinks]{hyperref}
\hypersetup{pdftitle={Fixed-window golden-ratio subset sums and Fibonacci representation counts}}

\newtheorem{theorem}{Theorem}[section]
\newtheorem{lemma}[theorem]{Lemma}
\newtheorem{proposition}[theorem]{Proposition}
\newtheorem{corollary}[theorem]{Corollary}
\newtheorem*{theoremA}{Theorem A}
\newtheorem*{theoremB}{Theorem B}
\newtheorem*{theoremC}{Theorem C}
\theoremstyle{definition}
\newtheorem{definition}[theorem]{Definition}
\newtheorem{example}[theorem]{Example}
\theoremstyle{remark}
\newtheorem{remark}[theorem]{Remark}

\numberwithin{equation}{section}
\raggedbottom

\newcommand{\Z}{\mathbb{Z}}
\newcommand{\R}{\mathbb{R}}
\newcommand{\val}{v}
\newcommand{\Rep}{\operatorname{Rep}}

\title[Fixed-window golden-ratio subset sums]{Fixed-window golden-ratio subset sums\\ and Fibonacci representation counts}

\author{HYDRA Research Project}
\thanks{Human-directed, AI-authored mathematical research. The proofs have not undergone independent human mathematical verification.}

\subjclass[2020]{Primary 11B39; Secondary 11A63, 05A17, 68R15}
\keywords{golden ratio; base phi; Fibonacci representations; Zeckendorf representation; subset sums; representation function}

\begin{document}

\begin{abstract}
Let $\varphi=(1+\sqrt5)/2$ and, for $m\ge 0$ and an integer $t$, let $\mu_m(t)$ be the number of
subsets $S\subseteq\{-2m,\dots,2m\}$ with $\sum_{k\in S}\varphi^k=t$. We prove that
$\mu_m(t)=R(F_{2m+2}\,t)$ for $0\le t\le L_{2m+1}$, where $R(n)$ is the number of representations of
$n$ as a sum of distinct Fibonacci numbers, and that $\mu_m(t)=0$ for all other integers $t$; the
correspondence is an explicit shift of exponents, and the proof uses only the algebraic conjugation
of $\mathbb{Z}[\varphi]$. Combining this transfer with a symmetry and a forced-digit argument, we
reduce the extremal problem for $\mu_m$ to results of Koc\'abov\'a, Mas\'akov\'a and Pelantov\'a on
the maxima of $R$ over intervals of fixed Zeckendorf length. We obtain
$\max_{t\in\mathbb{Z}}\mu_m(t)=F_{2m+1}$ for every $m\ge 0$, and we show that for $m\ge 2$ the
maximum is attained exactly at $t\in\{F_{2m},F_{2m+1},2F_{2m},F_{2m+2}\}$. We compare the transfer
with a theorem of Dekking and van Loon on natural base-$\varphi$ expansions.
\end{abstract}

\maketitle

\section{Introduction}\label{sec:intro}

Every positive integer has a unique finite expansion in base $\varphi=(1+\sqrt5)/2$ with digits
$0,1$ and no two adjacent ones (Bergman~\cite{Bergman1957}). Without the adjacency restriction an
integer has infinitely many finite expansions, because the golden mean flip $100\to 011$ (from
$\varphi^{n+1}=\varphi^n+\varphi^{n-1}$) can be applied indefinitely at the lowest digit.
Dekking and van Loon~\cite{DekkingVanLoon2024} obtained finite counts by two truncation rules
(Knott expansions and \emph{natural} expansions) and related the second one to the number $R(n)$
of representations of $n$ as a sum of distinct Fibonacci numbers, a function with a long history
\cite{Klarner1966,Robbins1996,BicknellFielder1999,Berstel2001,KMP2005,Stockmeyer2009,ChowSlattery2021}.

In this note we count expansions inside a \emph{fixed} symmetric window of exponents. For
$m\ge 0$ and $t\in\Z$ put
\[
  \mu_m(t)=\#\Bigl\{S\subseteq\{-2m,\dots,2m\}:\ \sum_{k\in S}\varphi^k=t\Bigr\},
  \qquad M_m=\max_{t\in\Z}\mu_m(t).
\]
Equivalently, $\mu_m(t)$ is the number of subsets of $\{\varphi^{-2m},\dots,\varphi^{2m}\}$ whose
sum is the integer $t$. Our first result identifies $\mu_m$ with a dilated restriction of $R$.

\begin{theoremA}[Theorem~\ref{thm:transfer}]
For every $m\ge0$ and every integer $t$ with $0\le t\le L_{2m+1}$,
\[
  \mu_m(t)=R(F_{2m+2}\,t),
\]
and $\mu_m(t)=0$ for all other integers $t$. The bijection behind the identity is the shift
$S\mapsto S+2m+2$ of exponents.
\end{theoremA}

The interval in Theorem~A is optimal: for every integer $t>L_{2m+1}$ the left-hand side vanishes
while $R(F_{2m+2}t)\ge1$. The essential point of the proof is that a shifted digit string must
reproduce \emph{both} coordinates of $t\varphi^{2m+2}=tF_{2m+2}\varphi+tF_{2m+1}$ in the basis
$\{\varphi,1\}$; matching the coefficient of $\varphi$ alone is not enough
(Remark~\ref{rem:onecoef}). We show that inside the stated range the second coordinate is forced by
a short conjugation estimate (Lemma~\ref{lem:conjbound}).

Theorem~A does not by itself determine $M_m$: the maximum of $R$ over all arguments below
$F_{4m+3}$ is $F_{2m+2}$, which is larger than the true value. Using a symmetry
$t\leftrightarrow L_{2m+1}-t$ and the observation that, after this symmetry, every representation has
two digits fixed, we reduce the problem to windows of $4m-1$ consecutive exponents, where the
results of Koc\'abov\'a, Mas\'akov\'a and Pelantov\'a~\cite{KMP2005} apply sharply.

\begin{theoremB}[Theorem~\ref{thm:max}]
For every $m\ge 0$, $M_m=F_{2m+1}$.
\end{theoremB}

\begin{theoremC}[Theorem~\ref{thm:argmax}]
For $m\ge2$, $\mu_m(t)=F_{2m+1}$ holds exactly for the four distinct integers
$t\in\{F_{2m},\,F_{2m+1},\,2F_{2m},\,F_{2m+2}\}$. For $m=0$ the maximizers are $0,1$, and for $m=1$
they are $1,2,3$.
\end{theoremC}

\subsection*{Relation to earlier work}
For integers $t$ whose Bergman expansion reaches the lowest exponent $-2m$ (these are exactly the
integers with $L_{2m-1}<t\le L_{2m+1}$), the window representations of $t$ are precisely its natural
expansions, and Theorem~A coincides with \cite[Theorem~4.4]{DekkingVanLoon2024}. For
$0\le t\le L_{2m-1}$ the window contains expansions that are not natural, and the multiplier
$F_{2m+2}$ in Theorem~A does not depend on $t$; we did not find this extension in the literature we
consulted (Section~\ref{sec:prior}). The extremal input for Theorems~B and~C is
\cite[Theorem~4.7, Lemma~3.1 and Corollary~5.1(i)]{KMP2005}; our contribution there is the reduction
from the window problem to the Fibonacci problem, including the transfer of the equality cases
with both coordinates in $\Z[\varphi]$. The lower bound in Theorem~B can also be read off from
\cite[Theorem~4.3]{DekkingVanLoon2024}; see Remark~\ref{rem:dvl43}. We make no claim of priority
beyond what is stated here.

\subsection*{Organization}
Section~\ref{sec:prelim} collects notation and the tools we use, including the translation of the
results of \cite{KMP2005} into the indexing of this paper. Section~\ref{sec:support} proves the
support and symmetry properties of $\mu_m$, Section~\ref{sec:transfer} proves Theorem~A,
Section~\ref{sec:max} proves Theorem~B, and Section~\ref{sec:argmax} proves Theorem~C.
Section~\ref{sec:examples} lists small cases and boundary examples, Section~\ref{sec:prior} compares
the results with the literature, and Section~\ref{sec:limits} records limitations and open questions.

\section{Notation and tools}\label{sec:prelim}

\subsection{Fibonacci and Lucas numbers}
Let $F_0=0$, $F_1=1$, $F_{n+2}=F_{n+1}+F_n$, and $L_0=2$, $L_1=1$, $L_{n+2}=L_{n+1}+L_n$. We extend
$F_n$ to negative indices by $F_{-n}=(-1)^{n+1}F_n$, so that the recurrence holds for all $n\in\Z$.
Let $\psi=1-\varphi=-\varphi^{-1}$. Binet's formulas are $F_n=(\varphi^n-\psi^n)/\sqrt5$ and
$L_n=\varphi^n+\psi^n$. For every $k\in\Z$,
\begin{equation}\label{eq:powers}
  \varphi^k=F_k\varphi+F_{k-1},\qquad \psi^k=F_k\psi+F_{k-1}.
\end{equation}
We use the following identities, valid for $m\ge0$ (and, in (d), for $m\ge1$):
\begin{align}
 &\text{(a)}\ \ \sum_{j=-m}^{m}\varphi^{2j}=L_{2m+1}; \qquad
  \text{(b)}\ \ F_{2m+2}L_{2m+1}=F_{4m+3}-1; \label{eq:idA}\\
 &\text{(c)}\ \ L_{2m+1}=F_{2m}+F_{2m+2}=F_{2m+1}+2F_{2m}; \nonumber\\
 &\text{(d)}\ \ F_{2m}F_{2m+2}=F_{2m+1}^2-1=\sum_{i=1}^{m}F_{4i}; \nonumber\\
 &\text{(e)}\ \ \varphi^{2m+1}=L_{2m+1}+\varphi^{-2m-1}. \nonumber
\end{align}
For (a), note that $\varphi^{2j+1}-\varphi^{2j-1}=\varphi^{2j}(\varphi-\varphi^{-1})=\varphi^{2j}$, so
the sum telescopes to $\varphi^{2m+1}-\varphi^{-2m-1}=\varphi^{2m+1}+\psi^{2m+1}$. Identity (b) is
the case $a=2m+1$, $b=2m+2$ of $L_aF_b=F_{a+b}+(-1)^aF_{b-a}$. Identity (c) is
$L_n=F_{n-1}+F_{n+1}$. In (d) the first equality is Cassini's identity, and the second follows by
induction from $F_{2m+2}F_{2m+4}-F_{2m}F_{2m+2}=F_{2m+2}L_{2m+2}=F_{4m+4}$. Identity (e) is
$L_{2m+1}=\varphi^{2m+1}+\psi^{2m+1}$. In particular $L_{2m+1}<\varphi^{2m+1}$, and an integer $t$
satisfies $t\le L_{2m+1}$ if and only if $t<\varphi^{2m+1}$.

Throughout, we write
\[
  C_m:=L_{2m+1}\qquad(m\ge0).
\]

\subsection{\texorpdfstring{The ring $\Z[\varphi]$}{The ring Z[phi]}}
Every element of $\Z[\varphi]$ is uniquely $a\varphi+b$ with $a,b\in\Z$, and it is a rational
number if and only if $a=0$. The conjugation $\sigma(a\varphi+b)=a\psi+b$ is a ring automorphism of
$\Z[\varphi]$ that fixes $\Z$. For a finite set $S\subseteq\Z$ we write
\[
  \val(S)=\sum_{k\in S}\varphi^k\in\Z[\varphi].
\]
By \eqref{eq:powers}, $\val(S)\in\Z$ if and only if $\sum_{k\in S}F_k=0$, and an equality
$\val(S)=\val(S')$ means equality of \emph{both} coordinates in the basis $\{\varphi,1\}$. Since
$\psi^k=(-1)^k\varphi^{-k}$, the term $\psi^k$ is positive for even $k$ and negative for odd $k$.

\subsection{Fibonacci representations}
For an integer $n\ge0$ let
\[
  \Rep(n)=\Bigl\{A\subseteq\{2,3,4,\dots\}\ \text{finite}:\ \sum_{j\in A}F_j=n\Bigr\},
  \qquad R(n)=\#\Rep(n).
\]
Thus $R(0)=1$ (the empty set), and the values $F_2,F_3,\dots=1,2,3,5,\dots$ are distinct, so $R$
counts representations of $n$ as sums of distinct Fibonacci numbers; this is sequence A000119
of~\cite{OEIS}. By Zeckendorf's theorem \cite{Lekkerkerker1952,Zeckendorf1972}, every $n\ge0$ has
exactly one representation $Z(n)\in\Rep(n)$ containing no two consecutive integers. In particular
$R(n)\ge1$ for all $n\ge 0$. If $A\in\Rep(n)$ and $n<F_{K+1}$, then $\max A\le K$.

\begin{lemma}[flip consistency]\label{lem:flip}
If $A,B\subseteq\{2,3,\dots\}$ are finite and $\sum_{j\in A}F_j=\sum_{j\in B}F_j$, then
$\val(A)=\val(B)$.
\end{lemma}

\begin{proof}
Let $n=\sum_{j\in A}F_j$. If $A$ contains two consecutive integers, let $j$ be the largest integer
with $j,j+1\in A$. Then $j+2\notin A$, since otherwise $j+1,j+2$ would be a higher consecutive pair.
Put $A'=(A\setminus\{j,j+1\})\cup\{j+2\}$. Because $F_{j+2}=F_{j+1}+F_j$ and
$\varphi^{j+2}=\varphi^{j+1}+\varphi^j$, the set $A'$ has the same Fibonacci sum and the same value
$\val(A')=\val(A)$, and it has one element fewer. Iterating, we reach after finitely many steps a set
without consecutive integers, which is $Z(n)$. Hence $\val(A)=\val(Z(n))$, and likewise
$\val(B)=\val(Z(n))$.
\end{proof}

\begin{lemma}[conjugate bound]\label{lem:conjbound}
For every finite set $A\subseteq\{2,3,\dots\}$,
\[
  -\varphi^{-2}<\sum_{j\in A}\psi^j<\varphi^{-1}.
\]
\end{lemma}

\begin{proof}
The positive terms come from even $j\ge2$; their sum is less than
$\sum_{i\ge1}\varphi^{-2i}=1/(\varphi^2-1)=\varphi^{-1}$. The negative terms come from odd $j\ge3$; their
sum exceeds $-\sum_{i\ge1}\varphi^{-2i-1}=-\varphi^{-1}\cdot\varphi^{-1}=-\varphi^{-2}$. Both
inequalities are strict because $A$ is finite.
\end{proof}

\subsection{\texorpdfstring{Maxima of $R$ on Zeckendorf-length intervals}{Maxima of R on Zeckendorf-length intervals}}
Koc\'abov\'a, Mas\'akov\'a and Pelantov\'a~\cite{KMP2005} index the Fibonacci numbers by
$F_0=F_1=1$; their $F_j$ is our $F_{j+1}$, and their function $R$ coincides with ours. Their
$\operatorname{Max}(\ell)$ is the maximum of $R$ over $[F^{\mathrm{KMP}}_\ell,F^{\mathrm{KMP}}_{\ell+1})
=[F_{\ell+1},F_{\ell+2})$, the set of integers whose Zeckendorf word has length $\ell$. A Zeckendorf
word is written $1\,0^{r_1}\,1\,0^{r_2}\cdots 1\,0^{r_l}$ from the most significant digit, so that
$r_l\ge0$ is the number of trailing zeros. In our indexing the following statements are
\cite[Theorem~4.7]{KMP2005}, \cite[Lemma~3.1(2)]{KMP2005} and \cite[Corollary~5.1(i)]{KMP2005}.

\begin{theorem}[\cite{KMP2005}]\label{thm:kmp}
For $\ell\ge1$ let $\mathcal M(\ell)=\max\{R(n):F_{\ell+1}\le n<F_{\ell+2}\}$.
\begin{enumerate}
\item[(a)] $\mathcal M(2h+1)=F_{h+2}$ for $h\ge0$, $\mathcal M(2h+2)=2F_{h+1}$ for $h\ge1$, and
$\mathcal M(2)=1$.
\item[(b)] For $k\ge1$, $\mathcal M(4k+3)$ is attained at exactly two integers. Their Zeckendorf
words have zero-block lengths $(r_1,\dots,r_l)=(1,3,\dots,3,4)$ with $k-1$ entries $3$, and
$(3,\dots,3,2)$ with $k$ entries $3$.
\end{enumerate}
\end{theorem}

For $m\ge1$ put
\[
  x_m=\sum_{i=1}^{m}F_{4i}=F_{2m}F_{2m+2},\qquad
  x_m'=\sum_{i=1}^{m-1}F_{4i+2}+F_{4m}.
\]
A word of length $4m-1=4(m-1)+3$ ends at our index $2$ and begins at our index $4m$. The word with
zero blocks $(3,\dots,3,2)$ has its ones at our indices $4,8,\dots,4m$, so it is $Z(x_m)$; the word
with zero blocks $(1,3,\dots,3,4)$ has its ones at $6,10,\dots,4m-2$ and $4m$, so it is $Z(x_m')$.
Summing even-indexed Fibonacci numbers gives $x_m+x_m'=F_{4m+2}-2$.

\begin{corollary}\label{cor:kmp}
Let $n=2q+1\ge1$.
\begin{enumerate}
\item[(a)] $R(u)\le F_{q+2}$ for every integer $0\le u<F_{n+2}$.
\item[(b)] If $m\ge2$, $0\le u<F_{4m+1}$ and $R(u)=F_{2m+1}$, then $u\in\{x_m,x_m'\}$.
\item[(c)] $R(x_m)=F_{2m+1}$ for every $m\ge1$.
\end{enumerate}
\end{corollary}

\begin{proof}
(a) We have $R(0)=1$. If $u\ge1$, then $F_{\ell+1}\le u<F_{\ell+2}$ for a unique $\ell$ with
$1\le\ell\le n$, and $R(u)\le\mathcal M(\ell)$. For odd $\ell=2h+1\le 2q+1$ this is
$F_{h+2}\le F_{q+2}$. For even $\ell=2h+2\le2q$ with $h\ge1$ it is $2F_{h+1}\le2F_q\le F_{q+2}$, and
$\mathcal M(2)=1$.

(b) Here $q=2m-1$ and $F_{q+2}=F_{2m+1}$. If $\ell\le 4m-2$, then the bounds in (a) give at most
$F_{2m}$ for odd $\ell$ and at most $2F_{2m-1}$ for even $\ell$; both are smaller than
$F_{2m+1}=F_{2m}+F_{2m-1}$ because $F_{2m-1}<F_{2m}$ for $m\ge2$. Also $R(0)=\mathcal M(2)=1$. Hence
$R(u)=F_{2m+1}$ forces $\ell=4m-1$, and Theorem~\ref{thm:kmp}(b) with $k=m-1\ge1$ leaves only the two
words described above, namely $Z(x_m)$ and $Z(x_m')$.

(c) For $m=1$, $x_1=3$ and $R(3)=2=F_3$. For $m\ge2$, \cite[Lemma~3.1(2)]{KMP2005} gives the value
$F_{2m+1}$ for the word with zero blocks $(3,\dots,3,2)$, which is $Z(x_m)$; this is also the
statement $\operatorname{Max}(4k+3)=\varrho(3,\dots,3,2)$ in \cite[Corollary~5.1(i)]{KMP2005}.
\end{proof}

\begin{remark}
The value $R(F_{2m+1}^2-1)=F_{2m+1}$ is also contained in a result of
Bicknell-Johnson~\cite{BicknellJohnson1999} on the smallest integer with a prescribed number of
representations, as noted in \cite[Remark after
Proposition~5.4]{KMP2005}, and it is an equality case of Stockmeyer's bound
$R(n)\le\sqrt{n+1}$~\cite{Stockmeyer2009}.
\end{remark}

\section{Support and symmetry}\label{sec:support}

\begin{lemma}[support]\label{lem:support}
Let $m\ge0$ and $S\subseteq\{-2m,\dots,2m\}$ with $\val(S)=t\in\Z$. Then $0\le t\le C_m$. If
moreover $m\ge1$ and $S\subseteq\{-2m+1,\dots,2m-1\}$, then $t\le C_{m-1}$.
\end{lemma}

\begin{proof}
All powers of $\varphi$ are positive, so $t\ge0$. Since $\sigma$ fixes integers,
$t=\sigma(t)=\sum_{k\in S}\psi^k$. Discarding the negative terms (odd $k$) and adding the missing
positive terms (even $k$ in the window), we obtain
$t\le\sum_{j=-m}^{m}\varphi^{-2j}=C_m$ by \eqref{eq:idA}(a). In the second case the even exponents
available are $-2m+2,\dots,2m-2$, which gives $t\le C_{m-1}$.
\end{proof}

\begin{lemma}[symmetry]\label{lem:sym}
For all $m\ge0$ and $t\in\Z$, $\mu_m(t)=\mu_m(C_m-t)$.
\end{lemma}

\begin{proof}
Let $E$ be the set of even integers in $[-2m,2m]$ and, for $S\subseteq[-2m,2m]$, put
$S^*=\{-k:k\in S\mathbin{\triangle}E\}$. The map $S\mapsto S^*$ is an involution of the subsets of
the window. Using $\varphi^{-k}=\psi^k$ for even $k$ and $\varphi^{-k}=-\psi^k$ for odd $k$,
\[
  \val(S^*)=\sum_{k\in E\setminus S}\psi^k-\sum_{k\in S\setminus E}\psi^k
  =\sum_{k\in E}\psi^k-\sum_{k\in S}\psi^k=\sigma(C_m)-\sigma(\val(S))=C_m-\sigma(\val(S)).
\]
If $\val(S)=t\in\Z$, then $\val(S^*)=C_m-t$.
\end{proof}

\section{The fixed-window transfer}\label{sec:transfer}

\begin{theorem}\label{thm:transfer}
Let $m\ge0$, $s=2m+2$, and let $t$ be an integer with $0\le t\le C_m$. Then $S\mapsto S+s$ is a
bijection from $\{S\subseteq\{-2m,\dots,2m\}:\val(S)=t\}$ onto $\Rep(F_s t)$. In particular,
\[
  \mu_m(t)=R(F_{2m+2}\,t)\qquad(0\le t\le C_m).
\]
\end{theorem}

\begin{proof}
\emph{The map is well defined and injective.} If $\val(S)=t$, then $S+s\subseteq\{2,\dots,4m+2\}$
and $\val(S+s)=t\varphi^s=tF_s\varphi+tF_{s-1}$. Comparing the coefficients of $\varphi$ by
\eqref{eq:powers} gives $\sum_{j\in S+s}F_j=tF_s$, so $S+s\in\Rep(F_st)$. Injectivity is clear.

\emph{The map is surjective.} Let $A\in\Rep(F_st)$. First, $F_{\max A}\le tF_s\le C_mF_s=F_{4m+3}-1$ by
\eqref{eq:idA}(b), so $A\subseteq\{2,\dots,4m+2\}$. Next, by \eqref{eq:powers},
$\val(A)=tF_s\varphi+b$ with $b=\sum_{j\in A}F_{j-1}\in\Z$. Hence
\[
  n:=\val(A)-t\varphi^s=b-tF_{s-1}\in\Z.
\]
Applying $\sigma$ gives $n=\sum_{j\in A}\psi^j-t\psi^s$. Since $s$ is even, $t\psi^s=t\varphi^{-s}$,
and since $t\le C_m<\varphi^{2m+1}$ by \eqref{eq:idA}(e),
\[
  0\le t\varphi^{-s}<\varphi^{2m+1}\varphi^{-2m-2}=\varphi^{-1}.
\]
Together with Lemma~\ref{lem:conjbound} this yields
$-1=-\varphi^{-2}-\varphi^{-1}<n<\varphi^{-1}<1$, so $n=0$ and $\val(A)=t\varphi^s$. The set
$S=A-s$ is contained in $\{-2m,\dots,2m\}$ and satisfies $\val(S)=\varphi^{-s}\val(A)=t$, and
$S+s=A$.
\end{proof}

\begin{corollary}\label{cor:support}
Let $m\ge0$.
\begin{enumerate}
\item[(a)] $\mu_m(t)>0$ if and only if $t\in\{0,1,\dots,C_m\}$.
\item[(b)] For an integer $t\ge0$, the identity $\mu_m(t)=R(F_{2m+2}t)$ holds if and only if
$t\le C_m$.
\end{enumerate}
\end{corollary}

\begin{proof}
By Lemma~\ref{lem:support}, $\mu_m(t)=0$ unless $0\le t\le C_m$. For $0\le t\le C_m$,
Theorem~\ref{thm:transfer} gives $\mu_m(t)=R(F_{2m+2}t)\ge1$. For $t>C_m$ we have $\mu_m(t)=0<1\le
R(F_{2m+2}t)$.
\end{proof}

\begin{remark}[one coefficient is not enough]\label{rem:onecoef}
Take $m=1$, $t=5>C_1=4$ and $S=\{-1,1,2\}$. Then $S+4=\{3,5,6\}$ and $F_3+F_5+F_6=15=5F_4$, but
$\val(S)=\varphi^{-1}+\varphi+\varphi^2=3\varphi$. The constant coordinates differ:
$F_2+F_4+F_5=9\neq 5F_3=10$; in the notation of the proof of Theorem~\ref{thm:transfer}, $n=-1$.
Thus the coefficient of $\varphi$ alone does not identify the value, and the hypothesis
$t\le C_m$ is exactly what forces the second coordinate.
\end{remark}

\begin{remark}[an alternative argument]\label{rem:altproof}
Surjectivity can also be derived from Lemma~\ref{lem:flip}. If $t$ has one representation $S_0$ in
the window (for instance, for $1\le t\le C_m$ its Bergman expansion lies in the window by
\cite[Proposition~5.1(b)]{DekkingVanLoon2024} and $\varphi^{2m+1}>C_m$), then every $A\in\Rep(F_st)$
satisfies $\val(A)=\val(S_0+s)=t\varphi^s$ by Lemma~\ref{lem:flip}. The proof given above avoids
this input and yields the existence of a window representation as a by-product.
\end{remark}

\begin{remark}[natural expansions]\label{rem:natural}
For an integer $t\ge1$ let $-\rho(t)$ be the lowest exponent of its Bergman expansion. Dekking and
van Loon call an expansion of $t$ \emph{natural} if it does not end in the digits $011$ and its
lowest exponent is $-\rho(t)$, and they prove $\mathrm{Tot}_\nu(t)=R(F_{\rho(t)+2}\,t)$ for $t>3$,
where $\mathrm{Tot}_\nu(t)$ is the number of natural expansions
\cite[Theorem~4.4]{DekkingVanLoon2024}. By \cite[Proposition~5.1(b)]{DekkingVanLoon2024},
$\rho(t)=2m$ exactly when $C_{m-1}<t\le C_m$. For such $t$ every expansion has lowest exponent at most
$-2m$ and highest exponent at most that of the Bergman expansion, because every expansion arises from
the Bergman expansion by golden mean flips \cite[Theorem~2.1]{DekkingVanLoon2024}. Hence a window
expansion of $t$ has lowest exponent exactly $-2m$; it cannot end in $011$, since undoing that final
flip would produce an expansion with a higher lowest exponent. Conversely, every natural expansion
lies in the window, since its highest exponent is at most that of the Bergman expansion, which is at
most $2m$ because $t<\varphi^{2m+1}$. So for $C_{m-1}<t\le C_m$ the window representations are exactly the natural
expansions and Theorem~\ref{thm:transfer} coincides with \cite[Theorem~4.4]{DekkingVanLoon2024}. For
$0\le t\le C_{m-1}$ the two statements differ. For example, for $m=2$ and $t=4$ we have
$\rho(4)=2$ and $\mathrm{Tot}_\nu(4)=R(12)=1$, whereas $\mu_2(4)=R(32)=4$, the four window
representations being
\[
  \{2,0,-2\},\quad \{2,0,-3,-4\},\quad \{2,-1,-2,-3,-4\},\quad \{1,0,-1,-2,-3,-4\}.
\]
\end{remark}

\section{The exact maximum}\label{sec:max}

For a finite set $W\subseteq\Z$ and $x\in\R$ let $N_W(x)=\#\{S\subseteq W:\val(S)=x\}$.

\begin{lemma}[fibres on odd windows]\label{lem:fibre}
Let $a\in\Z$, $n=2q+1\ge1$ and $W=\{a,a+1,\dots,a+n-1\}$. Then $N_W(x)\le F_{q+2}$ for every
$x\in\R$.
\end{lemma}

\begin{proof}
Multiplication by $\varphi^{2-a}$ maps the subsets of $W$ bijectively onto the subsets of
$W'=\{2,\dots,n+1\}$ and multiplies values by $\varphi^{2-a}$, so we may assume $W=W'$. If $N_W(x)=0$
there is nothing to prove. Otherwise $x\in\Z[\varphi]$, and all $S$ with $\val(S)=x$ have the same
coefficient $U=\sum_{j\in S}F_j$ of $\varphi$. The complement map $S\mapsto W\setminus S$ is a
bijection from the fibre of $x$ onto the fibre of $\val(W)-x$, and it replaces $U$ by $U_W-U$, where
$U_W=\sum_{j=2}^{n+1}F_j=F_{n+3}-2$. Replacing $x$ by $\val(W)-x$ if necessary, we may assume
$U\le U_W/2<F_{n+2}$. The identity map sends the fibre of $x$ injectively into $\Rep(U)$. Hence
$N_W(x)\le R(U)\le F_{q+2}$ by Corollary~\ref{cor:kmp}(a).
\end{proof}

\begin{theorem}\label{thm:upper}
For all $m\ge0$ and $t\in\Z$, $\mu_m(t)\le F_{2m+1}$.
\end{theorem}

\begin{proof}
For $m=0$ the only subsets are $\emptyset$ and $\{0\}$, so $\mu_0\le1=F_1$. Let $m\ge1$. By
Lemma~\ref{lem:support} we may assume $0\le t\le C_m$, and by Lemma~\ref{lem:sym} we may replace $t$
by $s=\min(t,C_m-t)$, so that $0\le s\le C_m/2$. Let $S$ be a window representation of $s$. Since
$C_m<\varphi^{2m+1}<2\varphi^{2m}$, we have $s<\varphi^{2m}$, so $2m\notin S$.

\emph{Case I: $s\le C_{m-1}$.} As $C_{m-1}<\varphi^{2m-1}$ by \eqref{eq:idA}(e), also $2m-1\notin S$.
Thus $\mu_m(s)=N_W(s)$ with $W=\{-2m,\dots,2m-2\}$, a window of $4m-1$ consecutive integers, and
Lemma~\ref{lem:fibre} with $q=2m-1$ gives $\mu_m(s)\le F_{2m+1}$.

\emph{Case II: $s>C_{m-1}$.} Then $-2m\in S$, for otherwise $S\subseteq\{-2m+1,\dots,2m-1\}$ and
$s\le C_{m-1}$ by Lemma~\ref{lem:support}. Hence $S\mapsto S\setminus\{-2m\}$ is a bijection from the
window representations of $s$ onto the subsets of $W=\{-2m+1,\dots,2m-1\}$ with value
$s-\varphi^{-2m}$, and Lemma~\ref{lem:fibre} again gives $\mu_m(s)\le F_{2m+1}$.
\end{proof}

\begin{proposition}\label{prop:attain}
For $m\ge1$,
\[
  \mu_m(F_{2m})=\mu_m(F_{2m+1})=\mu_m(2F_{2m})=\mu_m(F_{2m+2})=F_{2m+1}.
\]
\end{proposition}

\begin{proof}
Since $F_{2m}\le C_m$, Theorem~\ref{thm:transfer}, identity \eqref{eq:idA}(d) and
Corollary~\ref{cor:kmp}(c) give
\[
  \mu_m(F_{2m})=R(F_{2m}F_{2m+2})=R(x_m)=F_{2m+1}.
\]
As $C_m-F_{2m}=F_{2m+2}$ by \eqref{eq:idA}(c), Lemma~\ref{lem:sym} gives the same value at
$F_{2m+2}$.

Now let $s=F_{2m+1}$. We have $C_{m-1}=F_{2m-2}+F_{2m}<F_{2m+1}$ and $F_{2m+1}\le2F_{2m}$, so
$C_{m-1}<s\le C_m/2$. By Case~II of the proof of Theorem~\ref{thm:upper}, $\mu_m(s)=N_W(s-\varphi^{-2m})$
with $W=\{-2m+1,\dots,2m-1\}$. Multiplying by $\varphi^{2m+1}$ maps $W$ onto $\{2,\dots,4m\}$ and the
value $s-\varphi^{-2m}$ onto
\[
  y=F_{2m+1}\varphi^{2m+1}-\varphi=(F_{2m+1}^2-1)\varphi+F_{2m+1}F_{2m}.
\]
By \eqref{eq:idA}(d) this equals $F_{2m}F_{2m+2}\varphi+F_{2m}F_{2m+1}=F_{2m}\varphi^{2m+2}$, and
\[
  \val(Z(x_m))=\sum_{i=1}^{m}\varphi^{4i}=\varphi^4\,\frac{\varphi^{4m}-1}{\varphi^4-1}
  =\varphi^{2}\,\frac{\varphi^{4m}-1}{\sqrt5}=\varphi^{2m+2}F_{2m},
\]
using $\varphi^4-1=\sqrt5\,\varphi^2$ and Binet's formula. So the fibre of $y$ on $\{2,\dots,4m\}$
contains $Z(x_m)$. Since $x_m<F_{4m+1}$, every $A\in\Rep(x_m)$ lies in $\{2,\dots,4m\}$, and by
Lemma~\ref{lem:flip} it satisfies $\val(A)=\val(Z(x_m))=y$. Therefore the injection of the fibre of
$y$ into $\Rep(x_m)$ is a bijection, and $\mu_m(F_{2m+1})=R(x_m)=F_{2m+1}$. Finally
$C_m-F_{2m+1}=2F_{2m}$, and Lemma~\ref{lem:sym} gives the value at $2F_{2m}$.
\end{proof}

\begin{theorem}\label{thm:max}
For every $m\ge0$, $M_m=F_{2m+1}$.
\end{theorem}

\begin{proof}
For $m=0$, $\mu_0(0)=\mu_0(1)=1=F_1$. For $m\ge1$ combine Theorem~\ref{thm:upper} and
Proposition~\ref{prop:attain}.
\end{proof}

\begin{remark}[attainment via natural expansions]\label{rem:dvl43}
For $m\ge1$ the integers $F_{2m+1}$ and $F_{2m+2}$ lie in $(C_{m-1},C_m]$, so by
Remark~\ref{rem:natural} their window representations are their natural expansions, and
\cite[Theorem~4.3]{DekkingVanLoon2024} gives
$\mu_m(F_{2m+2})=\mathrm{Tot}_\nu(F_{2m+2})=F_{2m+1}$ and
$\mu_m(F_{2m+1})=\mathrm{Tot}_\nu(F_{2m+1})=F_{2m+1}$. This yields an independent route to the lower
bound in Theorem~\ref{thm:max}.
\end{remark}

\section{Classification of the maximizers}\label{sec:argmax}

For $m\ge2$ put
\[
  Q_m=\sum_{i=0}^{m-1}\varphi^{4i+2},\qquad T_m=\sum_{j=0}^{4m-2}\varphi^{j}.
\]
As in the proof of Proposition~\ref{prop:attain}, $Q_m=\varphi^{2m}F_{2m}$, and
$T_m=(\varphi^{4m-1}-1)/(\varphi-1)=\varphi^{4m}-\varphi$. Grouping the terms of $T_m-Q_m$ (the
exponents in $\{0,\dots,4m-2\}$ not congruent to $2$ modulo $4$) and using
$\varphi^{j}+\varphi^{j+1}=\varphi^{j+2}$ twice, we obtain
\begin{equation}\label{eq:TQ}
  T_m-Q_m=\sum_{i=0}^{m-1}\varphi^{4i+2}+\sum_{i=0}^{m-2}\varphi^{4i+3}
  =\sum_{i=1}^{m-1}\varphi^{4i}+\varphi^{4m-2}.
\end{equation}

\begin{lemma}[equality in Lemma~\ref{lem:fibre}]\label{lem:equality}
Let $m\ge2$, $W=\{0,1,\dots,4m-2\}$ and $x\in\R$ with $N_W(x)=F_{2m+1}$. Then
$x\in\{Q_m,\,T_m-Q_m\}$.
\end{lemma}

\begin{proof}
Follow the proof of Lemma~\ref{lem:fibre} with $n=4m-1$ after multiplying by $\varphi^2$, which maps
$W$ onto $\{2,\dots,4m\}$. Let $x'\in\{x,T_m-x\}$ be the value for which $U\le U_W/2$. Then
$F_{2m+1}=N_W(x')\le R(U)\le F_{2m+1}$ with $U<F_{4m+1}$, so $R(U)=F_{2m+1}$ and $U\in\{x_m,x_m'\}$ by
Corollary~\ref{cor:kmp}(b). Moreover, the injection of the fibre of $\varphi^2x'$ into $\Rep(U)$ is a
bijection between finite sets of equal size; in particular $Z(U)$ is the image of a subset in the
fibre, so $\varphi^2x'=\val(Z(U))$. Now $\val(Z(x_m))=\sum_{i=1}^m\varphi^{4i}=\varphi^2Q_m$, and by
\eqref{eq:TQ}, $\val(Z(x_m'))=\sum_{i=1}^{m-1}\varphi^{4i+2}+\varphi^{4m}=\varphi^2(T_m-Q_m)$. Hence
$x'\in\{Q_m,T_m-Q_m\}$, and since the complement interchanges these two values, also
$x\in\{Q_m,T_m-Q_m\}$.
\end{proof}

\begin{theorem}\label{thm:argmax}
Let $m\ge2$. Then $\mu_m(t)=F_{2m+1}$ if and only if $t\in\{F_{2m},F_{2m+1},2F_{2m},F_{2m+2}\}$, and
these four integers are distinct. For $m=0$ the maximizers are $0,1$; for $m=1$ they are $1,2,3$.
\end{theorem}

\begin{proof}
For $m=0$ see Theorem~\ref{thm:max}. For $m=1$, a direct enumeration of the $32$ subsets of
$\{-2,\dots,2\}$ gives $(\mu_1(0),\dots,\mu_1(4))=(1,2,2,2,1)$.

Let $m\ge2$. The four values are attained by Proposition~\ref{prop:attain}, and
$F_{2m}<F_{2m+1}<2F_{2m}<F_{2m+2}$ because $F_{2m-1}<F_{2m}$. Conversely, suppose
$\mu_m(t)=F_{2m+1}$ and let $s=\min(t,C_m-t)$, so that $\mu_m(s)=F_{2m+1}$ by
Lemma~\ref{lem:sym}. We use the two cases of the proof of Theorem~\ref{thm:upper}, together with
\eqref{eq:powers} and $F_{-n}=(-1)^{n+1}F_n$, which give
\[
  \varphi^{-2m}=-F_{2m}\varphi+F_{2m+1},\qquad \varphi^{1-2m}=F_{2m-1}\varphi-F_{2m},\qquad
  \varphi^{2-2m}=-F_{2m-2}\varphi+F_{2m-1}.
\]

\emph{Case I: $s\le C_{m-1}$.} The fibre of $s$ on $\{-2m,\dots,2m-2\}$ has $F_{2m+1}$ elements;
multiplying by $\varphi^{2m}$ and applying Lemma~\ref{lem:equality} gives
$s\in\{\varphi^{-2m}Q_m,\ \varphi^{-2m}(T_m-Q_m)\}$. Here $\varphi^{-2m}Q_m=F_{2m}$, while
\[
  \varphi^{-2m}(T_m-Q_m)=\varphi^{2m}-\varphi^{1-2m}-F_{2m}=F_{2m-2}\varphi+F_{2m-1}
\]
is irrational since $F_{2m-2}\ge1$. Hence $s=F_{2m}$.

\emph{Case II: $s>C_{m-1}$.} The fibre of $s-\varphi^{-2m}$ on $\{-2m+1,\dots,2m-1\}$ has $F_{2m+1}$
elements; multiplying by $\varphi^{2m-1}$ and applying Lemma~\ref{lem:equality} gives
$s\in\{\varphi^{-2m}+\varphi^{1-2m}Q_m,\ \varphi^{-2m}+\varphi^{1-2m}(T_m-Q_m)\}$. The first value is
$\varphi^{-2m}+F_{2m}\varphi=F_{2m+1}$. For the second,
\[
  \varphi^{-2m}+\varphi^{2m+1}-\varphi^{2-2m}-F_{2m}\varphi
  =(F_{2m+1}+F_{2m-2}-2F_{2m})\varphi+(F_{2m+1}+F_{2m}-F_{2m-1})=2F_{2m},
\]
because $F_{2m+1}+F_{2m-2}-2F_{2m}=F_{2m-1}+F_{2m-2}-F_{2m}=0$ and
$F_{2m+1}+F_{2m}-F_{2m-1}=F_{2m+1}+F_{2m-2}=2F_{2m}$. Since $C_m=F_{2m+1}+2F_{2m}$ and
$2F_{2m}>F_{2m+1}$, the value $2F_{2m}$ exceeds $C_m/2$, so $s=F_{2m+1}$.

Thus $s\in\{F_{2m},F_{2m+1}\}$, and $t\in\{s,C_m-s\}=\{F_{2m},F_{2m+1},F_{2m+2},2F_{2m}\}$.
\end{proof}

\begin{remark}
Theorem~\ref{thm:argmax} for $m\ge2$ uses the equality classification
\cite[Corollary~5.1(i)]{KMP2005}, which requires $k=m-1\ge1$. For $m=1$ the general formula gives
the multiset $\{F_2,F_3,2F_2,F_4\}=\{1,2,2,3\}$, which agrees with the three maximizers $1,2,3$.
\end{remark}

\section{Small cases and boundary examples}\label{sec:examples}

Table~\ref{tab:small} lists $\mu_m(t)$ for $m\le3$, and Table~\ref{tab:max} lists $M_m$ and the set
of maximizers for $m\le7$. Both tables were computed by exact enumeration with integer pairs
$(a,b)$ representing $a\varphi+b$; they agree with Theorems~\ref{thm:transfer}, \ref{thm:max}
and~\ref{thm:argmax}, and they are not used in any proof.

\begin{table}[ht]
\caption{Values of $\mu_m(t)$ for $0\le t\le C_m$.}\label{tab:small}
\centering
\begin{tabular}{@{}rrl@{}}
\toprule
$m$ & $C_m$ & $\mu_m(0),\mu_m(1),\dots,\mu_m(C_m)$\\
\midrule
$0$ & $1$ & $1,1$\\
$1$ & $4$ & $1,2,2,2,1$\\
$2$ & $11$ & $1,3,4,5,4,5,5,4,5,4,3,1$\\
$3$ & $29$ & $1,4,6,8,7,10,10,9,13,12,11,6,10,13,12,$\\
    &      & $12,13,10,6,11,12,13,9,10,10,7,8,6,4,1$\\
\bottomrule
\end{tabular}
\end{table}

\begin{table}[ht]
\caption{The maximum $M_m$ and the maximizers of $\mu_m$.}\label{tab:max}
\centering
\begin{tabular}{@{}rrl@{}}
\toprule
$m$ & $M_m$ & maximizers\\
\midrule
$0$ & $1$ & $0,1$\\
$1$ & $2$ & $1,2,3$\\
$2$ & $5$ & $3,5,6,8$\\
$3$ & $13$ & $8,13,16,21$\\
$4$ & $34$ & $21,34,42,55$\\
$5$ & $89$ & $55,89,110,144$\\
$6$ & $233$ & $144,233,288,377$\\
$7$ & $610$ & $377,610,754,987$\\
\bottomrule
\end{tabular}
\end{table}

\begin{example}
For $m=2$ we have $x_2=F_4+F_8=24$ and $x_2'=F_6+F_8=29$, and $R(24)=R(29)=5$. Theorem~\ref{thm:transfer}
gives $\mu_2(3)=R(3F_6)=R(24)=5$; in the proof of Proposition~\ref{prop:attain} the integer
$t=5=F_5$ falls into Case~II, with $U=F_5^2-1=24$. The values $t=6$ and $t=8$ are the reflections
$C_2-5$ and $C_2-3$.
\end{example}

\begin{example}[boundary]
For every $m$, $\mu_m(C_m)=\mu_m(0)=1=R(F_{2m+2}C_m)=R(F_{4m+3}-1)$, and $\mu_m(C_m+1)=0$ while
$R(F_{2m+2}(C_m+1))\ge1$. For $m=1$, Remark~\ref{rem:onecoef} exhibits a window subset whose shift is a
Fibonacci representation of $5F_4=15$ although its value is $3\varphi$.
\end{example}

\section{Comparison with earlier work}\label{sec:prior}

\emph{Dekking and van Loon~\cite{DekkingVanLoon2024}.} This is the closest source. Their
Theorem~4.4 is a transfer of the same type as Theorem~\ref{thm:transfer}, obtained by multiplying by
a power of $\varphi$; it counts natural expansions and uses the multiplier $F_{\rho(t)+2}$, which
depends on $t$. As explained in Remark~\ref{rem:natural}, the two statements coincide for
$C_{m-1}<t\le C_m$ and differ for $t\le C_{m-1}$, where the window also contains expansions that are
not natural. Their Theorem~4.3 gives the values at $F_{2m+1}$ and $F_{2m+2}$ used in
Remark~\ref{rem:dvl43}. Their paper does not study maxima of the counting functions.

\emph{Koc\'abov\'a, Mas\'akov\'a and Pelantov\'a~\cite{KMP2005}.} They determine the maximum of $R$ on
every interval $[F_{\ell+1},F_{\ell+2})$ and classify its arguments. Theorems~\ref{thm:max}
and~\ref{thm:argmax} use their Theorem~4.7, Lemma~3.1 and Corollary~5.1(i); without them our proofs
are incomplete. Their work does not involve powers of $\varphi$ or windows of exponents; the link is
provided by Lemma~\ref{lem:fibre}, by the forced digits in the proof of Theorem~\ref{thm:upper}, and by
Lemma~\ref{lem:equality}, which keeps track of both coordinates in $\Z[\varphi]$.

\emph{Other work on $R$.} Recursions and matrix formulas for $R$ go back to
\cite{Klarner1966,Robbins1996,BicknellFielder1999,Berstel2001}; Stockmeyer~\cite{Stockmeyer2009}
proved $R(n)\le\sqrt{n+1}$ with equality exactly at $n=F_k^2-1$, which contains
Corollary~\ref{cor:kmp}(c); Chow and Slattery~\cite{ChowSlattery2021} and Sanna~\cite{Sanna2025}
studied averages and power sums of $R$. None of these concerns windows of powers of $\varphi$.

\emph{Base-$\varphi$ representations and numeration systems.} Frougny and
Sakarovitch~\cite{FrougnySakarovitch1993} construct a finite automaton converting Fibonacci
representations into $\varphi$-expansions, and Frougny and Solomyak~\cite{FrougnySolomyak1992} study
finiteness and normalization of $\beta$-expansions. Normalization maps all representations of a
number to one normal form and therefore does not by itself give an equality of counts. Further
properties of $\varphi$-representations were proved with the Walnut theorem prover by
Shallit~\cite{Shallit2025} and by Shallit and Vukusic~\cite{ShallitVukusic2026}, and the structure of
the digits of base-$\varphi$ expansions is studied in \cite{Dekking2024Structure}; the extension of
\cite{KMP2005} to $m$-bonacci systems is \cite{KMP2007}.

Within the sources listed above we did not find Theorem~\ref{thm:transfer} for
$0\le t\le C_{m-1}$, Theorem~\ref{thm:max}, or Theorem~\ref{thm:argmax}. Our literature search did not
include a systematic search of MathSciNet or zbMATH, so this statement should be read as a report on
the sources consulted, not as a claim of priority.

\section{Limitations and further questions}\label{sec:limits}

The proofs of Theorems~\ref{thm:max} and~\ref{thm:argmax} rely on the published results
\cite[Theorem~4.7, Lemma~3.1, Corollary~5.1(i)]{KMP2005}, which we use as stated and do not reprove.
Theorem~\ref{thm:transfer} and Corollary~\ref{cor:support} are self-contained apart from Zeckendorf's
theorem. Exact computer checks confirm the statements of Theorems~\ref{thm:transfer},
\ref{thm:max} and~\ref{thm:argmax} for $m\le7$ by complete enumeration, the conclusion of
Theorem~\ref{thm:transfer} for the Zeckendorf representations of $F_{2m+2}t$ for all $t\le C_m$ and
$m\le12$, and Theorem~\ref{thm:kmp}(a) for all lengths $\ell\le34$ and Theorem~\ref{thm:kmp}(b) for
all lengths $\ell=4k+3\le31$; such checks do not replace the proofs. The following questions are natural.
\begin{enumerate}
\item A self-contained proof that $\max_{0\le u<F_{4m+1}}R(u)=F_{2m+1}$ with exactly the two
maximizers $x_m$ and $x_m'$ for $m\ge2$ would make Theorems~\ref{thm:max} and~\ref{thm:argmax}
independent of \cite{KMP2005}.
\item Lemma~\ref{lem:fibre} bounds every fibre, integral or not, on a window of odd length. Which
values of $x\in\Z[\varphi]$ attain the bound for general windows, and what is the maximum on windows of
even length?
\item Describe the second largest value of $\mu_m$ and the distribution of the values of $\mu_m$ on
$[0,C_m]$.
\item Asymmetric windows $\{a,\dots,b\}$, and analogous questions for other quadratic Pisot units,
are not treated here.
\end{enumerate}

\section*{Research provenance and AI contribution}
The results and proofs of this work were generated and audited by AI systems under human direction. The Owner's
public research-provenance classification is
\texttt{HUMAN\_DIRECTED\_AI\_AUTHORED}. The first complete proofs of the upper bound in
Theorem~\ref{thm:max}, of Theorem~\ref{thm:argmax} and of the transfer identity were
produced in HGS-E1-002 Request \\#14 by OpenAI model ``gpt-6-astra'' through the DeepSeek
Harness. This request-level identity is established; identities for other GOLD runs remain
unresolved. An independent reconstruction and adversarial check was carried out by Claude (Anthropic).
Literature audits were carried out by an AI system recorded as KIMI; an earlier numerical
verification was carried out through the DeepSeek API. The present text, including the
conjugation proof of Theorem~\ref{thm:transfer} and the attainment argument in
Proposition~\ref{prop:attain}, was drafted by Claude (historical model label
\texttt{claude-opus-5-5}). These contributor roles do not by themselves determine the
publication author list or the designation of individual systems as authors.

Independent human mathematical verification has not been completed. No human peer-review
status, Proof Verification certificate, or priority claim is asserted. The simplified proofs
in PUB-001 did not receive a separate independent audit. Formal submission metadata are maintained separately from this manuscript.

\bibliographystyle{amsplain}
\bibliography{references}

\end{document}
